% Watched literals report
% The table under tables/ is generated by tests/make_tables_watched.py from
% results/watched.csv, so `make bench-watched` refreshes it

\documentclass[11pt,a4paper]{article}
% Shared preamble for the NAIL094 reports
% Used by every task via % Shared preamble for the NAIL094 reports
% Used by every task via % Shared preamble for the NAIL094 reports
% Used by every task via \input{../../shared/preamble}

\usepackage[T1]{fontenc}
\usepackage[utf8]{inputenc}
\usepackage{lmodern}
\usepackage[margin=2.5cm]{geometry}

\usepackage{amsmath}
\usepackage{amssymb}
\usepackage{booktabs}
\usepackage{float}
\usepackage{graphicx}
\usepackage{listings}
\usepackage{xcolor}
\usepackage{caption}

\usepackage[hidelinks]{hyperref}

\setlength{\parskip}{0.5em}
\setlength{\parindent}{0pt}

% Title fields, set these in the report before \begin{document}
\makeatletter
\newcommand{\reportcourse}[1]{\def\@reportcourse{#1}}
\newcommand{\reporttask}[1]{\def\@reporttask{#1}}
\newcommand{\reportauthor}[1]{\def\@reportauthor{#1}}

\newcommand{\makereporttitle}{%
  \begin{center}
    {\LARGE\bfseries \@reportcourse \par}
    \vspace{0.8em}
    {\Large \@reporttask \par}
    \vspace{0.6em}
    {\normalsize \@reportauthor \par}
  \end{center}
  \vspace{1em}
  \hrule
  \vspace{1.5em}
}
\makeatother

% Inline code
\newcommand{\code}[1]{\texttt{#1}}

% Code listings
\lstdefinestyle{reportcode}{
  basicstyle=\ttfamily\small,
  keywordstyle=\color{blue!60!black},
  commentstyle=\color{gray},
  stringstyle=\color{green!40!black},
  breaklines=true,
  showstringspaces=false,
  frame=single,
  rulecolor=\color{gray!40},
  numbers=none,
}
\lstset{style=reportcode}


\usepackage[T1]{fontenc}
\usepackage[utf8]{inputenc}
\usepackage{lmodern}
\usepackage[margin=2.5cm]{geometry}

\usepackage{amsmath}
\usepackage{amssymb}
\usepackage{booktabs}
\usepackage{float}
\usepackage{graphicx}
\usepackage{listings}
\usepackage{xcolor}
\usepackage{caption}

\usepackage[hidelinks]{hyperref}

\setlength{\parskip}{0.5em}
\setlength{\parindent}{0pt}

% Title fields, set these in the report before \begin{document}
\makeatletter
\newcommand{\reportcourse}[1]{\def\@reportcourse{#1}}
\newcommand{\reporttask}[1]{\def\@reporttask{#1}}
\newcommand{\reportauthor}[1]{\def\@reportauthor{#1}}

\newcommand{\makereporttitle}{%
  \begin{center}
    {\LARGE\bfseries \@reportcourse \par}
    \vspace{0.8em}
    {\Large \@reporttask \par}
    \vspace{0.6em}
    {\normalsize \@reportauthor \par}
  \end{center}
  \vspace{1em}
  \hrule
  \vspace{1.5em}
}
\makeatother

% Inline code
\newcommand{\code}[1]{\texttt{#1}}

% Code listings
\lstdefinestyle{reportcode}{
  basicstyle=\ttfamily\small,
  keywordstyle=\color{blue!60!black},
  commentstyle=\color{gray},
  stringstyle=\color{green!40!black},
  breaklines=true,
  showstringspaces=false,
  frame=single,
  rulecolor=\color{gray!40},
  numbers=none,
}
\lstset{style=reportcode}


\usepackage[T1]{fontenc}
\usepackage[utf8]{inputenc}
\usepackage{lmodern}
\usepackage[margin=2.5cm]{geometry}

\usepackage{amsmath}
\usepackage{amssymb}
\usepackage{booktabs}
\usepackage{float}
\usepackage{graphicx}
\usepackage{listings}
\usepackage{xcolor}
\usepackage{caption}

\usepackage[hidelinks]{hyperref}

\setlength{\parskip}{0.5em}
\setlength{\parindent}{0pt}

% Title fields, set these in the report before \begin{document}
\makeatletter
\newcommand{\reportcourse}[1]{\def\@reportcourse{#1}}
\newcommand{\reporttask}[1]{\def\@reporttask{#1}}
\newcommand{\reportauthor}[1]{\def\@reportauthor{#1}}

\newcommand{\makereporttitle}{%
  \begin{center}
    {\LARGE\bfseries \@reportcourse \par}
    \vspace{0.8em}
    {\Large \@reporttask \par}
    \vspace{0.6em}
    {\normalsize \@reportauthor \par}
  \end{center}
  \vspace{1em}
  \hrule
  \vspace{1.5em}
}
\makeatother

% Inline code
\newcommand{\code}[1]{\texttt{#1}}

% Code listings
\lstdefinestyle{reportcode}{
  basicstyle=\ttfamily\small,
  keywordstyle=\color{blue!60!black},
  commentstyle=\color{gray},
  stringstyle=\color{green!40!black},
  breaklines=true,
  showstringspaces=false,
  frame=single,
  rulecolor=\color{gray!40},
  numbers=none,
}
\lstset{style=reportcode}


\reportcourse{NAIL094 --- Decision Procedures and SAT/SMT Solvers}
\reporttask{Watched Literals}
\reportauthor{Zerina Ahmetovic}

\begin{document}
\makereporttitle

\section{Task overview}

Lazy watched literals data structure is added to the DPLL solver of the
preceding task~\cite{dpllreport}.
Adjacency lists are to be used as an alternative and check
on the benchmark examples how much does watched literals save on
the number of checked clauses and running time.

Non-random benchmarks are to be used since random instances are not
good for comparison with adjacency lists (they do not have long clauses).
Therefore,the two non-random families from SATLIB~\cite{satlib} are
recommended for this task, pigeonhole~\cite{phole} and All Interval
Series~\cite{ais}, collected in Table~\ref{tab:comparison}. 
\section{Implementation}

As Watched Literals is an additon to the DPLL solver~\cite{dpllreport}, both of
them use
the same four-method interface and are chosen by a
template parameter of the solver, so the search loop, the trail and the decision
heuristic are the same code in both runs. Only propagation differs. The solver
is invoked as \texttt{dpll [input]} for adjacency lists and
\texttt{dpll --watched [input]} for watched literals.

\textbf{Adjacency lists.} For each literal $l$, \texttt{occ[l]} lists the
clauses containing it. When $l$ becomes true every clause holding $\lnot l$ has
lost a literal, so each of them is visited and rescanned from the start to see
whether it is satisfied, still open, unit, or falsified.

\begin{sloppypar}
\textbf{Watched literals.} Every clause of length at least two watches two of
its literals, and \texttt{watchList[l]} holds the clauses currently watching
$l$. A clause is inspected only when one of its two watched literals becomes
false; the engine then looks for any other non-false literal to move the watch
to, and only if none exists is the clause unit or falsified. Unit and empty
clauses cannot carry two watches and are handled once at initialisation.
\end{sloppypar}

The property that matters is what happens on backtracking: watches are not
restored, because a watch on a literal that has just become unassigned is still
a valid watch. Nothing is undone.

\textbf{Counting.} Both engines increment the same counter at the same point,
once per clause taken off the list being walked and immediately before the
clause is examined. This counts clause visits rather than distinct clauses, and
it is the only reading under which the two numbers are comparable.

\section{Experimental evaluation}

The times in Table~\ref{tab:comparison} were measured with g++ 14.2.0
(MSYS2 UCRT64) at \texttt{-std=c++20 -O2}. Each instance was run three times and
the fastest kept, the timer covers the search only. On every instance the two
engines report identical decision and conflict counts. They therefore walk the
same search tree, and the differences below are propagation and nothing else.

\begin{table}[H]
\centering
\small
\setlength{\tabcolsep}{4pt}
\begin{tabular}{lrrrlrrrrr}
\toprule
& & & & & \multicolumn{3}{c}{clauses checked} & \multicolumn{2}{c}{time (s)} \\
\cmidrule(lr){6-8}\cmidrule(lr){9-10}
instance & vars & clauses & $k>2$ & result & adjacency & watched & ratio & adjacency & watched \\
\midrule
hole6 & 42 & 133 & 0.05 & UNSAT & 63\,161 & 61\,730 & 1.02 & 0.00089 & 0.00128 \\
hole7 & 56 & 204 & 0.04 & UNSAT & 722\,098 & 707\,121 & 1.02 & 0.00939 & 0.01300 \\
hole8 & 72 & 297 & 0.03 & UNSAT & 9\,305\,155 & 9\,126\,435 & 1.02 & 0.12262 & 0.15621 \\
hole9 & 90 & 415 & 0.02 & UNSAT & 133\,426\,496 & 131\,029\,082 & 1.02 & 1.67565 & 2.21553 \\
ais6 & 61 & 441 & 0.28 & SAT & 2\,797 & 1\,855 & 1.51 & 0.00014 & 0.00012 \\
ais8 & 113 & 1149 & 0.27 & SAT & 51\,553 & 31\,531 & 1.64 & 0.00102 & 0.00084 \\
ais10 & 181 & 2377 & 0.26 & SAT & 1\,112\,800 & 653\,138 & 1.70 & 0.01909 & 0.01356 \\
ais12 & 265 & 4269 & 0.26 & SAT & 30\,580\,645 & 17\,250\,426 & 1.77 & 0.37316 & 0.33094 \\
\bottomrule
\end{tabular}
\caption{Adjacency lists against lazy watched literals on PHOLE and All Interval Series. $k>2$ is the fraction of clauses with more than two literals. Times are the fastest of three runs and cover search only. Both engines produce identical decision and conflict counts on every instance, so the two columns differ in propagation alone.}
\label{tab:comparison}
\end{table}


Table~\ref{tab:comparison} splits cleanly by family. On the pigeonhole instances
watched literals check only $2\%$ fewer clauses, a ratio of $1.02$ on all four,
and take $1.27$ to $1.45$ times as long to run. On All Interval Series they
check $34$ to $44\%$ fewer, a ratio rising from $1.51$ on \texttt{ais6} to
$1.77$ on \texttt{ais12}, and save $11$ to $29\%$ of the running time.

The $k>2$ column of Table~\ref{tab:comparison} accounts for the difference. A
watch only ever moves inside its own clause, so what the structure saves is
clause visits, and that is the quantity reported: adjacency lists
enter a clause whenever any of its $k$ literals is falsified, watched literals
only when one of the two being watched is. A binary clause has no third literal
for a watch to move to, and the pigeonhole instances are $95$ to $98\%$ binary:
there is nothing to save, and maintaining the watch lists is left as pure
overhead. All Interval Series carries $26$ to $28\%$ clauses longer than two
literals, which is enough for the lazy invariant to pay for itself. The same
reading explains the trend within each family --- the ratio grows with instance
size on All Interval Series, while on pigeonhole it stays at $1.02$ throughout.

\subsection*{Code}

\begin{sloppypar}
The full source is in the accompanying repository~\cite{repo}.
\texttt{src/watched\_literals.hpp} holds the watched engine and is the only file
added for this task; it includes \texttt{src/dpll.hpp} and nothing in
\texttt{dpll.hpp} refers back to it. The adjacency engine, the trail and the
search loop are unchanged in \texttt{src/dpll.hpp}.
\texttt{tests/check\_watched.py} asserts that the two engines agree on every
instance in the test suite and that their decision and conflict counts match,
and \texttt{tests/bench\_watched.py} produces the measurements.
\end{sloppypar}

\subsection*{Summary}

Watched literals reduce the number of clauses checked exactly where clauses are
long enough for a watch to have somewhere to move, and on those instances the
reduction carries through to the running time. On the pigeonhole family, which
is almost entirely binary clauses, they check barely fewer clauses and lose on
time to their own bookkeeping, as Table~\ref{tab:comparison} shows. What decides
the outcome is the clause width of the instance rather than its size or its
hardness.

\begin{thebibliography}{9}

\bibitem{dpllreport}
Z.~Ahmetovic.
\newblock DPLL experimental evaluation, the preceding report of this series.
\newblock Included in the accompanying repository.

\bibitem{repo}
Z.~Ahmetovic.
\newblock Source repository for the NAIL094 programming tasks.
\newblock \url{https://github.com/inferno73/SAT-solver-NAIL094}

\bibitem{satlib}
H.~H.~Hoos and T.~St\"utzle.
\newblock SATLIB: An Online Resource for Research on SAT.
\newblock \url{https://www.cs.ubc.ca/~hoos/SATLIB/benchm.html}

\bibitem{ais}
SATLIB benchmark problems, All Interval Series.
\newblock \url{https://www.cs.ubc.ca/~hoos/SATLIB/Benchmarks/SAT/AIS/descr.html}

\bibitem{phole}
SATLIB benchmark problems, DIMACS pigeon hole.
\newblock \url{https://www.cs.ubc.ca/~hoos/SATLIB/Benchmarks/SAT/DIMACS/PHOLE/descr.html}

\end{thebibliography}

\end{document}
